% =============================================================================
% Patch: §5 Toy anomaly and transverse obstruction
%   sec:toy-anomaly-and-transverse-obstruction
%
% Target file:   rh/rh_rebuild.tex
% Target lines:  1211 through 1263 (current §5 stub: bare definition + theorem
%                with two `wip' remarks).
%
% Action: this patch is NOT to be merged into rh_rebuild.tex yet.  It is
% staged under rh/patches/ for review.
%
% Migration content (relative to the current stub), restated in the
% rebuild's Gram-normalized real-phase framework:
%
%   - Definition def:off-line-toy-phase (replaces the bare stub
%     definition): explicit toy phase
%         Phi_toy(t; u, m_0, q_0) := q_0 (t - m_0)
%                                   + u^2 (1 - cos(q_0 (t - m_0))),
%     parameterized by the off-line displacement u = beta - 1/2, the
%     height m_0, and the local frequency q_0.  Real-valued and
%     C^infty in t for every (u, q_0).  The chosen u^2 modulation
%     amplitude (no extra q_0 factor) matches the natural q_0 scaling
%     of the §4 blocks G and N: q_+^{(2)} ~ q_0, q'^{(2)} ~ q_0^2,
%     q''^{(2)} ~ q_0^3.  After the natural diagonal scaling of
%     the §4 Gram blocks, the leading coefficient has a regular
%     high-frequency limit F_\infty(d) independent of q_0.  This is what
%     gives the visibility exponent A_toy = 0 on compact d-ranges avoiding
%     the zero set of F_\infty.
%
%   - Lemma lem:toy-phase-derivatives: explicit q_toy, q_toy', q_toy''
%     as elementary functions of (t - m_0), u, and q_0;
%     even-in-u^2 dependence and non-trivial first-order-in-(s/2)
%     splitting q_+ - q_- = 2 u^2 q_0 sin(q_0 s/2) at the symmetric
%     centers.
%
%   - Definition def:toy-finite-blocks: G_{toy,±}(s; m, u) and
%     N_toy(s; m, u) defined by applying the §4 package to Phi_toy.
%   - Definition def:toy-whitened-block: Omega_toy(s; m, u) by inverse
%     square root.
%
%   - Lemma lem:toy-block-positivity: G_{toy,±}(s) ≻ 0 on a retained
%     packet domain (compact d-range, |u| <= u_0).
%
%   - Definition def:toy-anomaly-functional: A_toy(X) := X_{12} + X_{21}
%     (the off-diagonal trace).  This is the rebuild's choice of
%     transverse functional, the analogue of v1's Phi_K
%     (proof_of_rh.tex line 1196).  A_toy(I_2) = 0;
%     A_toy(Omega_const) = 0 on the constant-phase on-line baseline by
%     the q_+ = q_- structure.
%
%   - Lemma lem:toy-anomaly-leading: leading-order expansion in u^2 of
%     the off-diagonal trace at small displacement,
%         A_toy(Omega_toy(s; m, u)) = u^2 . F_toy(d, q_0) + O(u^4),
%     with d := q_0 . s the dimensionless local separation.
%     F_toy(d, q_0) is a rational function of (sin d, cos d, sin(d/2),
%     cos(d/2), q_0) bounded uniformly on compact (d, q_0)-rectangles
%     and admits the high-frequency asymptotic limit
%         F_inf(d) := lim_{q_0 -> oo} F_toy(d, q_0)
%               = sqrt(3) (2 d^2 cos d + 5 d^2 cos 2d - 7 d^2
%                          + 6 d sin d - 15 d sin 2d
%                          - 12 cos 2d + 12) / (8 d^3 sin(d/2)),
%     which is q_0-independent, real-analytic on (0, infty), and
%     vanishes only at the discrete zero set of its trigonometric
%     numerator (with the sin(d/2) factor of the denominator
%     compensated by a matching factor in the numerator at d = 2 pi k).
%     On every compact subinterval D ⊂ (0, infty) avoiding the zeros
%     of F_inf there exist Q_0(D) >= 1 and c_toy(D) > 0 such that
%     |F_toy(d, q_0)| ≥ c_toy(D) for d in D and q_0 >= Q_0(D).
%
%   - Theorem thm:toy-anomaly-visibility (replaces the existing
%     theorem; same name and label): for every off-line displacement
%     u = beta - 1/2 with 0 < |u| < u_0(D) and every (m, s) with
%     d = q_0 s in a fixed compact D ⊂ (0, infty),
%         |A_toy(Omega_toy(s; m, u))| ≥ (1/2) c_toy(D) u^2.
%     Hence on the retained polynomial-weight subset of (m, s) on
%     which d ranges over D, the visibility exponent is
%     A_toy = 0 (no Q-dependence; the bound is c (beta - 1/2)^2 Q^0).
%
%   - Remark rem:toy-Pi-trans-deferred: the projection Pi_trans of
%     Section~\ref{sec:transverse-projection} is the further
%     gauge-fixing step.  Its definition is deferred to §6;
%     A_toy(.) = X_{12} + X_{21} of this section is one of the
%     functionals annihilating the value channel that Pi_trans is
%     designed to preserve, so the visibility statement transports to
%     the Pi_trans-projected form once §6 fixes the gauge.
%
% Closes (when migrated):
%   rem:wip-off-line-toy-phase                (closed: explicit formula).
%   rem:wip-toy-anomaly-visibility            (partially closed: leading
%                                              u^2 lower bound on compact
%                                              d-ranges; remainder still
%                                              flagged below).
%
% Remaining gaps after this patch:
%   - Uniform O(u^4) remainder bound across the full retained subset
%     (not just compact d-ranges avoiding F_toy-zeros).  Recorded as
%     rem:wip-toy-anomaly-uniform-remainder.
%   - The Pi_trans projection step itself is in §6.  This patch does
%     not migrate §6.
%   - The chosen toy phase is a model split-pair perturbation with
%     even-u^2 leading order, in line with v1's
%     "even symmetry of the off-line split" assumption (proof_of_rh.tex
%     line 51977).  A physically derived split-pair phase with the
%     same even-u^2 structure is not migrated here; the model toy
%     suffices for the visibility statement consumed downstream.
%
% Status (when merged):
%   \statusmig      not-started -> partial
%   \statusreferee  not-started -> in-progress
%   \statussympy    not-started -> verified
%       [rh/sympy/sec-toy-anomaly-and-transverse-obstruction.py]
%   \statussim      not-started -> verified
%       [rh/simulation/sec-toy-anomaly-and-transverse-obstruction.py]
%   \statuslean     not-started -> n/a
%
% Verification (every claim has a failing test in the cited script):
%   - rh/sympy/sec-toy-anomaly-and-transverse-obstruction.py
%       verify_toy_phase_derivatives,
%       verify_q_pm_split_at_centers,
%       verify_toy_blocks_substitution,
%       verify_omega_toy_baseline_at_u_zero,
%       verify_anomaly_functional_vanishes_on_baseline,
%       verify_leading_u2_expansion,
%       verify_F_toy_closed_form_q0_independent,
%       verify_F_toy_lower_bound_on_compact_d.
%   - rh/simulation/sec-toy-anomaly-and-transverse-obstruction.py
%       check_toy_phase_derivatives_mp,
%       check_toy_block_positivity_sweep,
%       check_anomaly_u2_scaling,
%       check_anomaly_lower_bound_on_compact_d,
%       check_anomaly_height_independence,
%       check_anomaly_sign_independence_in_u,
%       check_F_toy_against_sympy.
% =============================================================================

% =============================================================================
\section{Toy anomaly and transverse obstruction}
\label{sec:toy-anomaly-and-transverse-obstruction}
\statuspaper{LIVE}\quad
\statusmig{partial}\quad
\statusreferee{in-progress}\quad
\statussympy[rh/sympy/sec-toy-anomaly-and-transverse-obstruction.py]{verified}\quad
\statussim[rh/simulation/sec-toy-anomaly-and-transverse-obstruction.py]{verified}\quad
\statuslean{not-started}
\par\medskip

The toy phase \(\Ph_\toy\) defined below is a real-analytic visibility
model carrying an explicit even-\(u^2\) deformation of the linear
baseline, with \(u=\beta-\tfrac12\) the off-line displacement parameter.
It is not imported from the v1 toy matrix \(M(d)\); it replaces that
matrix-level toy package by a scalar-functional toy package for the
clean rebuild.  Applying the finite-block package of
Section~\ref{sec:finite-scale-local-blocks-and-whitening} to
\(\Ph_\toy\) produces a toy whitened block
\(\widehat\Om_\toy(s;m,u)\).  The off-diagonal trace
\(\mathcal A_\toy\) of this block is identically zero on the on-line
baseline and has nonzero leading-\(u^2\) coefficient
\(F_\toy(d,q_0)\), where \(d=q_0\,s\) is the dimensionless local
separation.  As \(q_0\to\infty\), this coefficient converges uniformly
on compact \(d\)-intervals to a \(q_0\)-independent closed-form limit
\(F_{\infty}(d)\).  The gauge-fixing projection
\(\Pi_\trans\) of Section~\ref{sec:transverse-projection} is a separate
step; this section proves the scalar visibility input that such a
projection must preserve.

\subsection{Off-line toy phase}
\label{subsec:off-line-toy-phase}

\begin{definition}[Off-line toy phase]
\label{def:off-line-toy-phase}
Fix a height \(m_0\), a local frequency \(q_0>0\), and an off-line
displacement \(u\in\mathbb R\).  The off-line toy phase is the
real-analytic phase
\begin{equation}
\label{eq:toy-phase}
\Ph_\toy(t;u,m_0,q_0)
\;:=\;q_0\,(t-m_0)\;+\;u^2\,\bigl(1-\cos(q_0\,(t-m_0))\bigr),
\qquad t\in\mathbb R.
\end{equation}
The displacement is identified with the off-line real part by
\(u=\beta-\tfrac12\).  By construction \(\Ph_\toy\) depends on \(u\)
only through \(u^2\) and is \(C^\infty\) on \(\mathbb R\) for every
\((u,q_0)\).
\end{definition}

\begin{lemma}[Toy phase derivatives]
\label{lem:toy-phase-derivatives}
With \(\tau:=t-m_0\), the derivatives of \(\Ph_\toy\) are
\begin{equation}
\label{eq:toy-phase-derivatives}
q_\toy(\tau)
=q_0+u^2\,q_0\sin(q_0\tau),
\quad
q_\toy'(\tau)=u^2\,q_0^{\,2}\cos(q_0\tau),
\quad
q_\toy''(\tau)=-u^2\,q_0^{\,3}\sin(q_0\tau).
\end{equation}
At \(\tau=0\),
\(q_\toy(0)=q_0\), \(q_\toy'(0)=u^2 q_0^{\,2}\),
\(q_\toy''(0)=0\).  The toy phase is not the retained theta chart of
Definition~\ref{def:local-phase-chart}; it is used only as a model
phase.  Its positivity analogue is
\[
        q_\toy(\tau)\ge q_0(1-u^2),
\]
so the toy same-point blocks are safely in the positive-frequency regime
whenever \(|u|<1\) and \(q_0(1-u^2)>0\).
\end{lemma}

\begin{proof}
Direct differentiation of \eqref{eq:toy-phase}.  The lower bound
\(q_\toy\ge q_0(1-u^2)\) follows from \(|\sin|\le 1\).
\end{proof}

\begin{lemma}[Symmetric-centre splitting at scale \(s\)]
\label{lem:q-pm-split-at-centres}
With \(t_\pm=m_0\pm s/2\) and \(d:=q_0\,s\), the toy phase derivatives
at the symmetric centres satisfy
\begin{equation}
\label{eq:q-pm-symmetric-split}
q_\pm:=q_\toy(\pm s/2)\;=\;q_0\;\pm\;u^2\,q_0\sin(d/2),
\qquad
q_\pm'\;=\;u^2\,q_0^{\,2}\cos(d/2),
\end{equation}
\[
q_\pm''\;=\;\mp\,u^2\,q_0^{\,3}\sin(d/2).
\]
In particular \(q_+ - q_- = 2 u^2 q_0\sin(d/2)\),
\(q_+' = q_-'\), and \(q_+''+q_-'' = 0\).  The phase difference
\(\Del_m^\toy(s)\) at the symmetric pair is
\begin{equation}
\label{eq:toy-Delta}
\Del_m^\toy(s)\;:=\;\Ph_\toy(t_-)-\Ph_\toy(t_+)\;=\;-q_0\,s\;=\;-d
\end{equation}
(no \(u\)-dependence: the cosine is even at \(\pm s/2\)).
\end{lemma}

\begin{proof}
\eqref{eq:q-pm-symmetric-split} substitutes \(\tau=\pm s/2\) into
\eqref{eq:toy-phase-derivatives} and uses oddness of \(\sin\) and
evenness of \(\cos\).  For \eqref{eq:toy-Delta}, evaluate
\eqref{eq:toy-phase} at \(\tau=\pm s/2\):
\[
\Ph_\toy(\pm s/2)
=\pm\tfrac{q_0 s}{2}
+u^2\bigl(1-\cos(\mp d/2)\bigr)
=\pm\tfrac{q_0 s}{2}
+u^2\bigl(1-\cos(d/2)\bigr),
\]
and subtract.
\end{proof}

\subsection{Toy whitened block}
\label{subsec:toy-whitened-block}

\begin{definition}[Toy finite blocks]
\label{def:toy-finite-blocks}
For \((m,s)\in\mathcal D_T\), with \(t_\pm=m\pm s/2\) and the toy phase
\(\Ph_\toy\) of Definition~\ref{def:off-line-toy-phase} centred at
\(m_0=m\), set
\begin{equation}
\label{eq:toy-G-N}
G_{\toy,\pm}(s;m,u)\;:=\;G_{m,\pm}(s)\big|_{\Ph=\Ph_\toy},
\qquad
N_\toy(s;m,u)\;:=\;N_m(s)\big|_{\Ph=\Ph_\toy},
\end{equation}
where the right-hand sides are the same-point and mixed blocks
\eqref{eq:Gpm}--\eqref{eq:Nm} of
Section~\ref{sec:finite-scale-local-blocks-and-whitening} evaluated
with \((q,q',q'',\Del_m)\) replaced by
\((q_\toy,q_\toy',q_\toy'',\Del_m^\toy)\) of
Lemmas~\ref{lem:toy-phase-derivatives}--\ref{lem:q-pm-split-at-centres}.
\end{definition}

\begin{definition}[Toy whitened block]
\label{def:toy-whitened-block}
For \((m,s)\in\mathcal D_T\) with \(G_{\toy,\pm}(s;m,u)\succ 0\),
\begin{equation}
\label{eq:omega-toy}
\widehat\Om_\toy(s;m,u)
\;:=\;G_{\toy,-}(s;m,u)^{-1/2}\,N_\toy(s;m,u)\,G_{\toy,+}(s;m,u)^{-1/2}.
\end{equation}
The inverse square root is the unique positive definite one.
\end{definition}

\begin{lemma}[Toy block positivity]
\label{lem:toy-block-positivity}
Fix a compact interval \(D=[d_-,d_+]\subset(0,\infty)\).  There exist
\(u_G=u_G(D)\in(0,\tfrac12]\), \(Q_G(D)\ge1\), and \(c_G(D)>0\) such
that whenever \(d=q_0\,s\in D\), \(|u|\le u_G\), and \(q_0\ge Q_G(D)\),
\[
        G_{\toy,\pm}(s;m,u)\succ 0,
        \qquad
        \lambda_{\min}\bigl(G_{\toy,\pm}(s;m,u)\bigr)\;\ge\;c_G(D)\,q_0.
\]
In particular \(\|G_{\toy,\pm}(s;m,u)^{-1/2}\|_{\op}\le C_G(D)q_0^{-1/2}\)
on the same region.
\end{lemma}

\begin{proof}
Write \(r=u^2\), \(z=\sin(d/2)\), and \(c=\cos(d/2)\).  From
\eqref{eq:q-pm-symmetric-split}, after factoring out the powers of
\(q_0\), the normalized same-point determinant has the form
\[
q_0^{-4}D_J(t_\pm)
=
4(1\pm rz)^4 \mp 2rz(1\pm rz)-3r^2c^2.
\]
This expression converges uniformly to \(4\) as \(r\to0\), uniformly in
\(d\in D\).  Decrease \(u_G\), if necessary, so that
\(q_0^{-4}D_J(t_\pm)\ge 1\) and \(q_\pm\ge q_0/2\) for
\(|u|\le u_G\).  By Lemma~\ref{lem:algebraic-same-point-gram-criterion},
\(G_{\toy,\pm}\succ0\).  Also, on \(d\in D\) and \(|u|\le u_G\),
\(\Tr G_{\toy,\pm}\le C(D)q_0^3\).  Therefore
\[
\lambda_{\min}(G_{\toy,\pm})
\ge \frac{\det G_{\toy,\pm}}{\Tr G_{\toy,\pm}}
\ge c_G(D)q_0.
\]
After increasing \(Q_G(D)\), this holds throughout the retained
high-frequency region.  The inverse-square-root bound follows by taking
reciprocals and square roots.
\end{proof}

\subsection{Toy anomaly functional and leading expansion}
\label{subsec:toy-anomaly-functional}

\begin{definition}[Toy anomaly functional]
\label{def:toy-anomaly-functional}
For a real \(2\times 2\) matrix \(X=(x_{ij})\), the toy anomaly
functional is the off-diagonal trace
\begin{equation}
\label{eq:toy-anomaly-functional}
\mathcal A_\toy(X)\;:=\;x_{12}+x_{21}.
\end{equation}
\end{definition}

\begin{lemma}[Vanishing on the on-line baseline]
\label{lem:A-toy-baseline-vanishing}
At \(u=0\) the toy phase reduces to the linear baseline
\(\Ph_\toy(t;0,m,q_0)=q_0(t-m)\), and \(\widehat\Om_\toy(s;m,0)\) is
the constant-phase whitened block of the linear baseline.  In
particular
\begin{equation}
\label{eq:A-toy-baseline}
        \mathcal A_\toy\bigl(\widehat\Om_\toy(s;m,0)\bigr)\;=\;0
        \qquad\bigl((m,s)\in\mathcal D_T^\times\bigr).
\end{equation}
\end{lemma}

\begin{proof}
At \(u=0\), \eqref{eq:toy-phase-derivatives} gives \(q_\pm=q_0\) and
\(q_\pm'=q_\pm''=0\); by \eqref{eq:Gpm},
\(G_{\toy,\pm}(s;m,0)=(1/\pi)\diag(2q_0,q_0^{\,3}/6)\) is diagonal,
and \(G_{\toy,-}^{-1/2}=G_{\toy,+}^{-1/2}\) is the same diagonal
matrix.  By \eqref{eq:Nm} with \(q_+=q_-=q_0\),
\[
N_\toy(s;m,0)_{12}+N_\toy(s;m,0)_{21}
=\frac{(q_+-q_-)\cos\Del_m^\toy(s)}{\pi s}=0.
\]
Whitening by the diagonal \(G^{-1/2}\) preserves this off-diagonal
cancellation: \(\widehat\Om_{12}+\widehat\Om_{21}
=(G^{-1/2})_{11}(G^{-1/2})_{22}(N_{12}+N_{21})=0\).
\end{proof}

\begin{lemma}[Leading \(u^2\) expansion of the off-diagonal trace]
\label{lem:toy-anomaly-leading}
Fix a compact interval \(D=[d_-,d_+]\subset(0,\infty)\).  There exist
\(Q_1=Q_1(D)\ge1\), \(u_1=u_1(D)>0\), and a function
\(F_\toy:D\times[Q_1,\infty)\to\mathbb R\) such that for every
\(0<|u|<u_1\), every \(q_0\ge Q_1\), and every
\((m,s)\in\mathcal D_T^\times\) with \(d=q_0\,s\in D\),
\begin{equation}
\label{eq:toy-anomaly-leading}
\mathcal A_\toy\bigl(\widehat\Om_\toy(s;m,u)\bigr)
\;=\;u^2\,F_\toy(d,q_0)\;+\;u^4\,R_\toy(d,q_0,u),
\end{equation}
with \(|R_\toy(d,q_0,u)|\le C_R(D)\) uniformly on the region
\(\{|u|\le u_1,\,d\in D,\,q_0\ge Q_1(D)\}\).  The high-frequency limit
\begin{equation}
\label{eq:F-toy-asymptotic}
        F_{\infty}(d)\;:=\;\lim_{q_0\to\infty}F_\toy(d,q_0)
\end{equation}
exists, is reached uniformly on compact \(d\)-intervals, and admits
the explicit closed form
\begin{equation}
\label{eq:F-toy-closed-form}
\begin{aligned}
F_{\infty}(d)
\;&=\;\frac{\sqrt 3}{8\,d^{\,3}\,\sin(d/2)}\;\Bigl(
2\,d^{\,2}\cos d+5\,d^{\,2}\cos 2d-7\,d^{\,2}\\
&\qquad\qquad\qquad\quad
+6\,d\sin d-15\,d\sin 2d-12\cos 2d+12\Bigr).
\end{aligned}
\end{equation}
\(F_{\infty}\) is real-analytic on \((0,\infty)\); the apparent
\(\sin(d/2)\)-pole at \(d=2\pi k\) is removable, and \(F_{\infty}\)
vanishes only on the discrete zero set of its trigonometric
numerator.  On every compact subinterval \(D\subset(0,\infty)\)
avoiding the zeros of \(F_{\infty}\) there exist \(Q_0(D)\ge 1\) and
\(c_\toy(D)>0\) with
\begin{equation}
\label{eq:F-toy-lower-bound}
        \bigl|F_\toy(d,q_0)\bigr|\;\ge\;c_\toy(D)
        \qquad (d\in D,\;q_0\ge Q_0(D)).
\end{equation}
\end{lemma}

\begin{proof}
Each entry of \(G_{\toy,\pm}(s;m,u)\) and \(N_\toy(s;m,u)\) is a
real-analytic function of \(r:=u^2\) by \eqref{eq:toy-phase-derivatives}
and \eqref{eq:Gpm}--\eqref{eq:Nm}.  To make the uniformity in \(q_0\)
explicit after substituting \(s=d/q_0\), use the natural powers
\(q_0,q_0^2,q_0^3\) attached to the value/derivative coordinates:
\[
G_{11}=q_0\,g_{11},\quad G_{12}=q_0^2\,g_{12},\quad
G_{22}=q_0^3\,g_{22},
\]
and similarly
\[
N_{11}=q_0\,n_{11},\quad N_{12}=q_0^2\,n_{12},\quad
N_{21}=q_0^2\,n_{21},\quad N_{22}=q_0^3\,n_{22}.
\]
The coefficients \(g_{ij}\) and \(n_{ij}\) are real-analytic functions
of \((r,d,\eta)\), \(\eta=q_0^{-1}\), on
\([0,u_1^2]\times D\times[0,\eta_1]\), once \(u_1\) is small and
\(Q_1=\eta_1^{-1}\) is large.  By Lemma~\ref{lem:toy-block-positivity},
the normalized determinants of the \(G_{\toy,\pm}\) stay bounded away
from zero on this compact parameter set.  The explicit \(2\times2\)
positive-definite inverse-square-root formula therefore gives
\[
(G_{\toy,\pm}^{-1/2})_{11}=q_0^{-1/2}h_{11}^{\pm},\quad
(G_{\toy,\pm}^{-1/2})_{12}=q_0^{-3/2}h_{12}^{\pm},\quad
(G_{\toy,\pm}^{-1/2})_{22}=q_0^{-3/2}h_{22}^{\pm},
\]
where the \(h_{ij}^{\pm}\) are again real-analytic and uniformly bounded
on the same compact set.  Multiplying the three factors in
\(G_{\toy,-}^{-1/2}N_\toy G_{\toy,+}^{-1/2}\) shows that every entry of
\(\widehat\Om_\toy\), and in particular
\(\mathcal A_\toy(\widehat\Om_\toy)\), is real-analytic in \(r\) with
uniformly bounded Taylor coefficients.  Since the \(u^0\) coefficient
vanishes by Lemma~\ref{lem:A-toy-baseline-vanishing}, the leading
nonzero term is at order \(u^2\); call its coefficient
\(F_\toy(d,q_0)\).

Direct computation using the closed-form \(2\times 2\) SPD
inverse-square-root, applied to
\(G_{\toy,\pm}=G_{\toy,\pm}^{(0)}+u^2\,G_{\toy,\pm}^{(2)}+O(u^4)\) of
\eqref{eq:Gpm} with the toy substitutions of
\eqref{eq:q-pm-symmetric-split} and to \(N_\toy\) of
Definition~\ref{def:toy-finite-blocks}, yields an explicit expression
for \(F_\toy(d,q_0)\) in
\((\sin d,\cos d,\sin(d/2),\cos(d/2),q_0)\).  After choosing
\(Q_1(D)\) large enough, this expression is regular for
\(d\in D\), \(q_0\ge Q_1(D)\), has a finite \(q_0\to\infty\) limit
uniformly on \(D\), and the limiting expression is
\eqref{eq:F-toy-closed-form}.

The factor \(\sin(d/2)\) of the denominator of
\eqref{eq:F-toy-closed-form} vanishes at \(d=2\pi k\),
\(k\in\mathbb Z_{\ge 1}\); the trigonometric numerator vanishes at
the same points.  Expanding the numerator at each \(d=2\pi k\) shows a
factor of \(\sin(d/2)\), so the apparent poles are removable and
\(F_{\infty}\) extends real-analytically to \((0,\infty)\).  On any compact subinterval avoiding the zeros of
\(F_{\infty}\), the difference \(F_\toy(d,q_0)-F_{\infty}(d)\) is
\(O(1/q_0)\) uniformly; for \(q_0\) sufficiently large
\(|F_\toy(d,q_0)|\ge \tfrac12|F_{\infty}(d)|\), giving
\eqref{eq:F-toy-lower-bound}.  The remainder \(R_\toy\) is the \(u^4\)-and-higher part of this
real-analytic expansion.  The compactness of the normalized parameter
set and the uniform positive-definiteness just described give a
Cauchy-type bound for these coefficients, uniformly for \(d\in D\),
\(|u|\le u_1\), and \(q_0\ge Q_1(D)\), after decreasing \(u_1\) and
increasing \(Q_1(D)\) if necessary.
\end{proof}

\subsection{Toy anomaly visibility}
\label{subsec:toy-anomaly-visibility}

\begin{theorem}[Toy anomaly visibility]
\label{thm:toy-anomaly-visibility}
Let \(D\subset(0,\infty)\) be a compact interval avoiding the
isolated zeros of \(F_{\infty}\) of
Lemma~\ref{lem:toy-anomaly-leading}.  There are constants
\(Q_\toy(D)\ge1\) and \(u_\toy(D)>0\) such that, for every off-line
displacement \(u=\beta-\tfrac12\in\mathbb R\) with \(0<|u|\le u_\toy(D)\),
every \(q_0\ge Q_\toy(D)\), and every \((m,s)\in\mathcal D_T^\times\)
with \(d=q_0\,s\in D\), the toy
whitened block \(\widehat\Om_\toy(s;m,u)\) of
Definition~\ref{def:toy-whitened-block} satisfies
\begin{equation}
\label{eq:toy-anomaly-visibility-A}
        \bigl|\mathcal A_\toy\bigl(\widehat\Om_\toy(s;m,u)\bigr)\bigr|
        \;\ge\;
        \tfrac12\,c_\toy(D)\,u^2.
\end{equation}
The constant \(c_\toy(D)>0\) is the lower bound of
Lemma~\ref{lem:toy-anomaly-leading} and depends on \(D\) only.
On the polynomial-weight subset of midpoints and separations on which
\(d\in D\), the visibility exponent is \(A_\toy=0\): the lower bound
is \((c_\toy(D)/2)\,u^2\,Q^0\), independent of the height.
\end{theorem}

\begin{proof}
Choose \(Q_\toy(D)\ge\max\{Q_1(D),Q_G(D)\}\).  Let \(u_G(D)\) be the
smallness threshold from Lemma~\ref{lem:toy-block-positivity}.  Choose
\[
        u_\toy(D)\le \min\{u_1(D),u_G(D)\}
\]
small enough that both lemmas apply.  By
\eqref{eq:toy-anomaly-leading} and \eqref{eq:F-toy-lower-bound},
\[
\bigl|\mathcal A_\toy\bigl(\widehat\Om_\toy(s;m,u)\bigr)\bigr|
\;\ge\;
u^2\bigl(c_\toy(D)-u^2\,C_R(D)\bigr).
\]
Decreasing \(u_\toy(D)\) so that \(u_\toy(D)^{\,2}\,C_R(D)\le c_\toy(D)/2\) gives
\eqref{eq:toy-anomaly-visibility-A}.
\end{proof}

\begin{remark}[Transverse projection deferred]
\label{rem:toy-Pi-trans-deferred}
The functional \(\mathcal A_\toy\) annihilates the on-line baseline
(Lemma~\ref{lem:A-toy-baseline-vanishing}).  The later projection
\(\Pi_\trans\) must be chosen so that this scalar functional descends to
or is represented on the retained transverse quotient.  Once
Section~\ref{sec:transverse-projection} proves that compatibility, the
visibility bound \eqref{eq:toy-anomaly-visibility-A} transfers to the
\(\Pi_\trans\)-projected comparison with the same exponent
\(A_\toy=0\).  This compatibility is not proved in the present section.
\end{remark}


\begin{remark}[Compatibility with the archived v1 toy matrix]
\label{rem:toy-phase-v1-compatibility}
The toy phase \eqref{eq:toy-phase} is a replacement toy for the clean
scalar-functional rebuild.  It is not the archived v1 matrix \(M(d)\),
and this section does not import the v1 calibration map
\(\Psi_{x,d}\) or the Frobenius pairings with \(A_{\val}\).  Consequently
this change is isolated for the present rebuild through
Sections~\ref{sec:transverse-projection}--\ref{sec:projective-rigidity},
where only the scalar functional \(\mathcal A_\toy\), the exponent
\(A_\toy=0\), and the lower bound \eqref{eq:toy-anomaly-visibility-A}
are consumed.  Any later attempt to import a v1 theorem that explicitly
uses \(M(d)\), \(\Phi_K(M(d))\), \(A_{\val}\), or \(\Psi_{x,d}\) must be
rederived for the present toy phase or kept in the archive.
\end{remark}

\begin{wip}\label{rem:wip-toy-anomaly-uniform-remainder}
The bound \eqref{eq:toy-anomaly-visibility-A} is stated on compact
intervals \(D\subset(0,\infty)\) avoiding the isolated zeros of
\(F_{\infty}\); a uniform extension to the entire retained subset
(allowing \(d\to 0\), \(d\to\infty\), or \(d\) approaching a
zero along the polynomial-weight sample) requires a sharper analysis
of the \(O(u^4)\) remainder \(R_\toy\) of
Lemma~\ref{lem:toy-anomaly-leading}, together with a polynomial
density estimate on the avoidance set of the \(F_\toy\)-zeros.
The numerical and symbolic verifications cover the compact case;
the uniform extension is deferred.
\end{wip}
